\documentclass[11pt,a4paper]{article}
\usepackage[spanish]{babel}
\usepackage[utf8]{inputenc}
\usepackage[T1]{fontenc}
\usepackage{amsmath,amssymb,bm}
\usepackage{geometry}
\usepackage{xcolor}
\usepackage{booktabs}
\usepackage{hyperref}
\geometry{margin=2.5cm}
\hypersetup{colorlinks=true,linkcolor=blue,urlcolor=blue}

\newcommand{\IS}{\delta}
\newcommand{\QS}{\Delta E_Q}
\newcommand{\BHF}{B_{\mathrm{HF}}}
\newcommand{\dd}{\mathrm{d}}

\title{Distribuciones de desplazamiento isomérico $P(\IS)$ y distribuciones $P(\IS,\QS)$ en espectros Mössbauer}
\author{Fitbauer -- nota matemática de implementación}
\date{\today}

\begin{document}
\maketitle

\section{Motivación física}

El desplazamiento isomérico $\IS$ mide el desplazamiento del centro del patrón
Mössbauer respecto a una referencia, normalmente $\alpha$-Fe. Está relacionado
con la densidad electrónica en el núcleo y, por tanto, con el estado de
oxidación, coordinación y covalencia del átomo de Fe. En materiales ordenados se
suele ajustar un único $\IS$ por fase. Sin embargo, en vidrios, silicatos,
materiales amorfos, ferritinas, nanopartículas o mezclas de entornos locales,
puede existir una distribución continua de desplazamientos.

Fitbauer implementa dos extensiones:
\begin{itemize}
  \item una distribución unidimensional $P(\IS)$;
  \item distribuciones bidimensionales que combinan $\IS$ con otra magnitud,
  especialmente $P(\IS,\QS)$ y $P(\BHF,\IS)$.
\end{itemize}
El caso $P(\IS,\QS)$ es particularmente útil para dobletes paramagnéticos anchos,
donde el centro del doblete y su separación pueden variar simultáneamente.

\section{Distribución unidimensional $P(\IS)$}

Para una malla de desplazamientos $\IS_i$, el modelo 1D es
\begin{equation}
  T(v_k)=b_0+b_1v_k-\sum_i P_i\,S(v_k;\IS_i,\QS,\BHF,\Gamma),
\end{equation}
donde $S$ es el patrón de absorción de profundidad unitaria. Si $\BHF=0$, el
patrón se reduce en la práctica a un doblete o singlete según el valor de $\QS$.
Los pesos cumplen
\begin{equation}
  P_i\ge0.
\end{equation}
El problema regularizado se escribe como
\begin{equation}
  \Phi(\bm p)=\|W(\bm y-X\bm p)\|^2+\alpha\|D^2\bm p\|^2,
\end{equation}
donde $D^2$ es el operador de segundas diferencias. La probabilidad normalizada
se calcula como
\begin{equation}
  p(\IS_i)=\frac{P_i}{\int P(\IS)\,\dd\IS}.
\end{equation}

\section{Distribución bidimensional $P(\IS,\QS)$}

Para una malla $\IS_i$, $Q_j$, el modelo es
\begin{equation}
  T(v_k)=b_0+b_1v_k-\sum_{ij}P_{ij}\,
  S(v_k;\IS_i,Q_j,\BHF,\Gamma).
\end{equation}
Si se desea describir una población paramagnética, se fija $\BHF=0$ y cada punto
$(\IS_i,Q_j)$ genera un doblete centrado en $\IS_i$ con separación $Q_j$:
\begin{align}
  v_{1,ij}&=\IS_i-\frac{Q_j}{2},\\
  v_{2,ij}&=\IS_i+\frac{Q_j}{2}.
\end{align}

La regularización 2D penaliza curvatura en las dos direcciones:
\begin{equation}
  \Phi(\bm P)=\|W(\bm y-X\bm z)\|^2
  +\alpha_{\IS}\|D_{\IS}^2P\|^2
  +\alpha_Q\|D_Q^2P\|^2,
  \qquad P_{ij}\ge0.
\end{equation}
En forma aumentada:
\begin{equation}
  \begin{bmatrix}
    WX\\
    \sqrt{\alpha_{\IS}}\,D_{\IS}^2\\
    \sqrt{\alpha_Q}\,D_Q^2
  \end{bmatrix}\bm z\simeq
  \begin{bmatrix}
    W\bm y\\
    \bm 0\\
    \bm 0
  \end{bmatrix}.
\end{equation}

\section{Interpretación de $P(\IS,\QS)$}

Las marginales son
\begin{align}
  P_{\IS}(\IS_i)&=\int P(\IS_i,Q)\,\dd Q,\\
  P_Q(Q_j)&=\int P(\IS,Q_j)\,\dd\IS.
\end{align}
Una nube inclinada en el mapa puede indicar que los sitios con mayor $\IS$ tienen
también mayor o menor $\QS$. Esto puede sugerir familias de entornos locales,
por ejemplo variación de valencia, coordinación o distorsión del poliedro.

Fitbauer calcula diagnósticos descriptivos del mapa:
\begin{align}
  \langle\IS\rangle,\quad \sigma_{\IS},\quad
  \langle Q\rangle,\quad \sigma_Q,\quad
  \rho_{\IS Q}.
\end{align}
La correlación aparente es
\begin{equation}
  \rho_{\IS Q}=\frac{\langle(\IS-\langle\IS\rangle)(Q-\langle Q\rangle)\rangle}
  {\sigma_{\IS}\sigma_Q}.
\end{equation}
Estos valores describen la distribución ajustada, pero no prueban por sí solos
que la distribución sea única o física.

\section{Distribución $P(\BHF,\IS)$}

También se implementa $P(\BHF,\IS)$:
\begin{equation}
  T(v_k)=b_0+b_1v_k-\sum_{ij}P_{ij}\,
  S(v_k;B_i,\QS,\IS_j,\Gamma).
\end{equation}
Puede ser útil en fases magnéticas donde el campo hiperfino y el desplazamiento
isomérico varían conjuntamente, por ejemplo por distribución de tamaños,
oxidación parcial o entornos de superficie. No obstante, es especialmente
sensible a errores de calibración de velocidad: un desplazamiento global de la
escala puede parecer una distribución de $\IS$.

\section{Riesgos de identificabilidad}

\textbf{Advertencia importante:} las distribuciones que contienen $\IS$ son muy
sensibles a calibración, folding y fondo. En un doblete,
\begin{equation}
  v_1=\IS-\frac{\QS}{2},\qquad v_2=\IS+\frac{\QS}{2},
\end{equation}
por lo que cambios de $\IS$ desplazan el centro y cambios de $\QS$ modifican la
separación. Con líneas anchas o ruido, muchas combinaciones $(\IS,\QS)$ pueden
producir espectros parecidos. Además, $\IS$ se correlaciona con:
\begin{itemize}
  \item el folding point;
  \item la calibración $V_{\max}$;
  \item la referencia isomérica;
  \item la anchura $\Gamma$;
  \item la pendiente/fondo;
  \item componentes discretas omitidas.
\end{itemize}
Por tanto, una distribución $P(\IS)$ o $P(\IS,\QS)$ visualmente suave y con buen
residuo puede carecer de significado físico si está compensando un error
instrumental o un modelo incompleto.

\section{Buenas prácticas}

Antes de interpretar una distribución con $\IS$ se recomienda:
\begin{enumerate}
  \item fijar o verificar cuidadosamente la calibración de velocidad;
  \item usar una referencia isomérica conocida;
  \item comprobar el folding point con el diagnóstico de residuo;
  \item comparar con modelos discretos y con $P(\QS)$ 1D;
  \item variar $\alpha_{\IS}$ y $\alpha_Q$ en varios órdenes de magnitud;
  \item repetir con distinto número de bins;
  \item evitar publicar detalles finos que desaparecen al cambiar la
  regularización;
  \item usar información externa: química, XRD, TEM, magnetometría o literatura.
\end{enumerate}
Para publicación, conviene mostrar que el mapa $P(\IS,\QS)$ es estable y que un
modelo más simple no explica los datos de forma satisfactoria.

\section{Estado en Fitbauer}

En la GUI se accede mediante los modos:
\begin{itemize}
  \item \textbf{P(IS)} para la distribución 1D;
  \item \textbf{P(IS,$\Delta E_Q$) 2D};
  \item \textbf{P(BHF,IS) 2D}.
\end{itemize}
Los modos 2D comparten el motor regularizado usado para $P(\BHF,\QS)$. La GUI
muestra mapa de calor, marginales, L-surface de regularización, exportación TSV,
heatmap interactivo Plotly e informe PDF con mapa 2D.

\end{document}
